\section{Results and Discussion}
\label{sec:results}

\subsection{Fault-type recovery: severity vs type}
Representations built on gas concentrations agree only weakly with the Duval fault types
(KMeans ARI~$0.16$; a plain autoencoder $0.14$) because they encode gassing \emph{magnitude}:
the total gas is linearly readable from the autoencoder's code ($R^2(m\mid\mathbf{z})=0.63$).
Replacing concentrations with the CLR composition removes magnitude by construction and
recovers fault \emph{type}: KMeans on the standardised 5-D CLR reaches ARI~$0.49$, and the
deployed two-dimensional PCA projection \textbf{ARI~$0.545\pm0.002$} (NMI~$0.54$), versus
$0.12$ for raw proportions (Fig.~\ref{fig:repr}a). The improvement is therefore due
specifically to the log-ratio geometry---not to normalisation, and \emph{not to the neural
network}: the SD-CAE autoencoder reaches only $0.47\pm0.06$, and the adversarial independence
term degrades it further (ARI $0.30$--$0.31$ at $\lambda=1$--$10$), confirming that part of the
composition--severity correlation is legitimate fault signal. We therefore deploy the linear
variant and report both neural components as negative ablations. The structure generalises
across time---fitting on pre-2022 samples and assigning 2022+ samples to the frozen centroids
yields ARI $0.45$---and is stable for zero-replacement $\delta\le10^{-3}$ (ARI $0.46$--$0.55$),
though it degrades at $\delta=10^{-2}$, an unavoidable sensitivity when acetylene is zero in
$97\%$ of samples; at unit level (latest sample per unit) agreement is $0.35$. The
two-dimensional code is a learned, label-free diagnostic map on which Duval classes separate
visibly (Fig.~\ref{fig:repr}b). Because Duval is a deterministic rule rather than an
attention-driven label, this is our most defensible result; the moderate absolute ARI reflects
that Duval itself is a heuristic over an imbalanced fleet, so we report the large
\emph{relative} gain.

\begin{figure*}[!t]
  \centering
  \includegraphics[width=0.9\textwidth]{representation_panel}
  \caption{Severity-disentangled compositional representation. (a) Fault-type agreement with
  Duval across representations (mean$\pm$std, five seeds): the gain comes from the CLR
  log-ratio geometry, and a linear PCA projection beats the SD-CAE autoencoder (negative
  ablation). (b) The deployed label-free diagnostic map---PCA-2 of the standardised
  CLR---coloured by Duval class; discharge (D1/D2), partial discharge, and thermal (T1--T3)
  regions separate.}
  \label{fig:repr}
\end{figure*}

\subsection{Label-free health/risk screening}
The risk index ranks the fleet with AUC~$0.74$ (unit-bootstrap $95\%$ CI $[0.68,0.80]$) against
\texttt{fault\_note}, versus $0.52$--$0.56$ for the conventional severity proxies---though a
trivial sample-count predictor reaches a statistically indistinguishable $0.76$ on this
confounded label (Fig.~\ref{fig:confound}a), so the comparison must be read against that floor
(Sec.~\ref{subsec:confound}). Its top decile carries a
$\sim$$2.4\times$ event rate and its top ten units a $\sim$$4.4\times$ rate; the event rate
falls from $\sim$27\% in the riskiest decile to $\sim$2\% in the safest, with the five riskiest
deciles above the fleet base rate (Fig.~\ref{fig:health}). A component ablation shows
the order of importance: severity ($0.64$) $\rightarrow$ $+$acetylene ($0.71$) $\rightarrow$
$+$temporal H2 rate ($0.73$) $\rightarrow$ $+$anomaly ($0.74$); acetylene-awareness is the
dominant lever, consistent with arcing being the most dangerous fault. Among anomaly detectors,
Isolation Forest ($0.60$) and autoencoder reconstruction ($0.58$) are weak but additive, whereas
Deep~SVDD collapses on this small tabular feature space (AUC $0.52$); a battery of nine standard
detectors yields no improvement over the deployed combination. The component set was chosen
during development with feedback from this same label, so the absolute AUC carries some design
optimism; the component ordering, not the third decimal, is the finding.

\begin{figure}[t]
  \centering
  \includegraphics[width=0.92\columnwidth]{health_decile_note_rate}
  \caption{Field-event rate by health/risk decile (D1 = riskiest). The rate falls from
  $\sim$27\% (riskiest) to $\sim$2\% (safest); the five riskiest deciles exceed the fleet base
  rate (dashed).}
  \label{fig:health}
\end{figure}

\begin{figure*}[!t]
  \centering
  \includegraphics[width=0.93\textwidth]{confound_panel}
  \caption{The surveillance-bias analysis. (a) A trivial sample count matches the physics index
  against the field-note label and exceeds every conventional severity baseline (the confound
  floor). (b) Strict temporal hold-out: the score's forward validity vanishes once the
  sample-count confound is controlled, while the count alone forward-predicts. (c) Confound-free
  chemistry target (arcing onset): gas chemistry predicts, the sampling confound does not.}
  \label{fig:confound}
\end{figure*}

\subsection{The field-event label is surveillance-confounded}
\label{subsec:confound}
This is the paper's cautionary result. A trivial predictor---how many times a unit has been
sampled---scores AUC~$\mathbf{0.76}$ against \texttt{fault\_note}, statistically
indistinguishable from the physics index ($0.74$; paired-bootstrap difference CI
$[-0.10,+0.06]$, $p=0.58$) and far above the conventional baselines (Fig.~\ref{fig:confound}a).
The mechanism is operational: utilities re-sample units they are worried about, so the sample
count encodes attention rather than gas chemistry, and $69\%$ of the notes contain Thai text
largely missed by the English keyword label. In a joint logistic model the count adds signal
far beyond the physics (likelihood-ratio $p=9\times10^{-13}$) while the physics retains a
small significant increment ($p=0.003$): the label is dominated by, though not purely, operator
attention. Repairing the label with our Thai-note translation (six additional own-unit
electrical events) makes the count's advantage disappear entirely (AUC $0.734$ vs $0.735$),
indicating that part of its edge came from label incompleteness. The consequence is decisive
under a strict temporal hold-out: a label-free score built from each unit's first-half history
predicts second-half events at AUC $0.54$, the first-half sample count alone predicts them
better ($0.68$), and once the score is residualised on that count its forward AUC falls to
chance ($0.50$, Fig.~\ref{fig:confound}b). Because the first-half sample count is known
\emph{before} the second-half events, this is a prospective effect, not reverse causation. We
therefore do not claim genuine forward fault prediction from \texttt{fault\_note}; any AUC
against it must be read against the $0.76$ confound floor.

To evaluate without this confound we use a chemistry-defined target---the onset of arcing
(acetylene appearing within two years in a unit that has none), which is uncorrelated with the
sample count (correlation $-0.01$ with the target). Across 3{,}794 prediction points the
positives are rare and clustered: 45 positive points arising from only 17 onset units.
Accounting for this clustering with a unit-blocked bootstrap, ethylene and hydrogen forecast
onset at AUC $0.64$ ($95\%$ CI $[0.52,0.77]$) and $0.63$ ($[0.53,0.73]$)---both excluding
chance---while the sample count is uninformative ($0.49$, CI $[0.38,0.58]$,
Fig.~\ref{fig:confound}c); the result is robust across nine onset-threshold and horizon variants
(AUC $0.64$--$0.73$). Thermal activity (ethylene) and incipient hydrogen thus precede
arcing---a physically sensible, confound-free early signal, modest in effect size and with wide
intervals owing to the rarity of the event, but distinguishable from chance.
